\documentclass[11pt,a4paper]{article}

\usepackage[utf8]{inputenc}
\usepackage[T1]{fontenc}
\usepackage{geometry}
\usepackage{graphicx}
\usepackage{booktabs}
\usepackage{array}
\usepackage{longtable}
\usepackage{listings}
\usepackage{xcolor}
\usepackage{hyperref}
\usepackage{fancyhdr}
\usepackage{titlesec}
\usepackage{amsmath}

\geometry{margin=2.5cm}

\hypersetup{
    colorlinks=true,
    linkcolor=blue!70!black,
    urlcolor=blue!70!black
}

% Code listing style
\definecolor{codebg}{gray}{0.95}
\definecolor{codegreen}{rgb}{0,0.5,0}
\definecolor{codeblue}{rgb}{0,0,0.7}

\lstdefinestyle{verilog}{
    backgroundcolor=\color{codebg},
    basicstyle=\ttfamily\footnotesize,
    keywordstyle=\color{codeblue}\bfseries,
    commentstyle=\color{codegreen}\itshape,
    stringstyle=\color{red!70!black},
    numbers=left,
    numberstyle=\tiny\color{gray},
    numbersep=8pt,
    frame=single,
    frameround=tttt,
    breaklines=true,
    captionpos=b,
    morekeywords={module,endmodule,input,output,wire,reg,assign,always,begin,end,
                  case,endcase,default,genvar,generate,endgenerate,for,parameter},
    tabsize=4
}

\lstset{style=verilog}

% Header/Footer
\pagestyle{fancy}
\fancyhf{}
\fancyhead[L]{\footnotesize 7-Bit Hybrid DAC --- Digital Definition}
\fancyhead[R]{\footnotesize CONFIDENTIAL}
\fancyfoot[C]{\thepage}
\renewcommand{\headrulewidth}{0.4pt}

% Title
\title{
    \vspace{-1cm}
    \textbf{Digital Implementation Definition} \\
    \Large 7-Bit Hybrid DAC Controller \\
    (4-Bit Binary + 3\textrightarrow 7 Thermometric with Reorder Mux)
}
\author{Retym --- Analog/Mixed-Signal Design Group}
\date{Document Date: \today \\ Revision: 1.1}

\begin{document}

\maketitle
\thispagestyle{fancy}

\tableofcontents
\newpage

%==============================================================================
\section{Overview}

\begin{table}[h]
\centering
\begin{tabular}{@{}ll@{}}
\toprule
\textbf{Parameter}           & \textbf{Value} \\
\midrule
Total DAC resolution         & 7 bits \\
Binary LSB segment           & 4 bits (codes 0--15) \\
Thermometric MSB segment     & 3 MSB $\rightarrow$ 7 thermometer-coded bits \\
MSB input coding             & 2's complement (range $-4$ to $+3$) \\
Total analog outputs         & 4 binary + 7 thermometric = 11 unit elements \\
\bottomrule
\end{tabular}
\caption{DAC Parameters}
\end{table}

\textbf{Input:} 7-bit digital code \texttt{D[6:0]}
\begin{itemize}
    \item \texttt{D[3:0]} --- Binary LSB segment (directly drives 4 binary-weighted current/voltage sources).
    \item \texttt{D[6:4]} --- 3-bit MSB segment in \textbf{2's complement} format. Converted to unsigned before thermometric encoding.
\end{itemize}

%==============================================================================
\section{Architecture Block Diagram}

\begin{verbatim}
                 +-------------------------------------------------------------------+
                 |                    DAC Digital Controller                          |
                 |                                                                    |
 D[6:0] ------->|  +----------+    +----------+    +------------------+              |
                 |  | 2's Comp |    | Binary-to|    | Thermometric     |              |
         D[6:4]-+->| to       +--->| Thermo   +--->| Reorder Mux      +---> T[6:0]  |
                 |  | Unsigned |    | Encoder  |    | (7x7)            |              |
                 |  +----------+    +----------+    +------------------+              |
                 |                                                                    |
         D[3:0]-+------------------------------------------------------> B[3:0]     |
                 |                                                                    |
                 +-------------------------------------------------------------------+

 T[6:0] -> 7 thermometric unit elements (MSB segment)
 B[3:0] -> 4 binary-weighted elements (LSB segment)
\end{verbatim}

%==============================================================================
\section{2's Complement to Unsigned Converter}

\subsection{Function}

Converts the 3-bit 2's complement MSB input \texttt{D[6:4]} (range $-4$ to $+3$) to a 3-bit unsigned value (range $0$ to $7$) suitable for the binary-to-thermometric encoder.

The conversion adds an offset of 4 (i.e., inverts the MSB):
\[
  \text{unsigned} = \text{twos\_comp} + 4 = \text{twos\_comp} \oplus \{1, 0, 0\}
\]
Equivalently: invert bit [2] (the sign bit) to produce the unsigned representation.

\subsection{Truth Table}

\begin{table}[h]
\centering
\begin{tabular}{@{}ccc@{}}
\toprule
\textbf{D[6:4] (2's complement)} & \textbf{Decimal value} & \textbf{Unsigned output} \\
\midrule
100 & $-4$ & 000 (0) \\
101 & $-3$ & 001 (1) \\
110 & $-2$ & 010 (2) \\
111 & $-1$ & 011 (3) \\
000 & $+0$ & 100 (4) \\
001 & $+1$ & 101 (5) \\
010 & $+2$ & 110 (6) \\
011 & $+3$ & 111 (7) \\
\bottomrule
\end{tabular}
\caption{2's Complement to Unsigned Conversion}
\end{table}

\subsection{Implementation}

\begin{lstlisting}[caption={2's Complement to Unsigned Converter}]
module twos_comp_to_unsigned (
    input  wire [2:0] twos_in,    // D[6:4], 2's complement
    output wire [2:0] unsigned_out // unsigned magnitude for thermo encoder
);

    // Invert MSB to convert from 2's complement to offset-binary (unsigned)
    assign unsigned_out = {~twos_in[2], twos_in[1:0]};

endmodule
\end{lstlisting}

\subsection{Timing}

\begin{itemize}
    \item Single inverter on MSB; effectively zero gate delay addition.
    \item No pipeline stage required.
\end{itemize}

%==============================================================================
\section{Binary-to-Thermometric Encoder}

\subsection{Function}

Converts 3-bit unsigned input (output of the 2's complement converter) to 7-bit thermometric output.

\subsection{Truth Table}

\begin{table}[h]
\centering
\begin{tabular}{@{}cc@{}}
\toprule
\textbf{Unsigned Input (decimal)} & \textbf{Thermometric Output TH[6:0]} \\
\midrule
000 (0) & 0000000 \\
001 (1) & 0000001 \\
010 (2) & 0000011 \\
011 (3) & 0000111 \\
100 (4) & 0001111 \\
101 (5) & 0011111 \\
110 (6) & 0111111 \\
111 (7) & 1111111 \\
\bottomrule
\end{tabular}
\caption{Binary-to-Thermometric Encoding}
\end{table}

\subsection{Implementation}

Combinational logic. RTL implementation:

\begin{lstlisting}[caption={Binary-to-Thermometric Encoder}]
module bin2thermo (
    input  wire [2:0] bin,   // unsigned value from converter
    output reg  [6:0] thermo // thermometric code
);

    always @(*) begin
        case (bin)
            3'd0: thermo = 7'b000_0000;
            3'd1: thermo = 7'b000_0001;
            3'd2: thermo = 7'b000_0011;
            3'd3: thermo = 7'b000_0111;
            3'd4: thermo = 7'b000_1111;
            3'd5: thermo = 7'b001_1111;
            3'd6: thermo = 7'b011_1111;
            3'd7: thermo = 7'b111_1111;
            default: thermo = 7'b000_0000;
        endcase
    end

endmodule
\end{lstlisting}

\subsection{Timing}

\begin{itemize}
    \item Pure combinational; latency = propagation delay of case decode ($\sim$1--2 gate delays).
    \item If pipeline stage needed for timing closure, register output with system clock.
\end{itemize}

%==============================================================================
\section{Thermometric Bit Reorder Mux}

\subsection{Purpose}

Enable arbitrary reordering of the 7 thermometric output bits. Use cases:
\begin{itemize}
    \item Compensate for layout-induced mismatch / gradient effects.
    \item DEM (Dynamic Element Matching) scrambling patterns.
    \item Post-silicon trimming / calibration.
    \item Debug / characterization flexibility.
\end{itemize}

\subsection{Functional Description}

A $7\times7$ crossbar mux. Each output position \texttt{T[i]} selects from any of the 7 thermometric encoder outputs \texttt{TH[6:0]} based on a per-output select field.

\begin{verbatim}
         TH[6:0] (from encoder)
           | | | | | | |
           v v v v v v v
     +-------------------------+
     |   7x7 Crossbar Mux     |
     |                         |
     |  sel[0] -> T[0]        |
     |  sel[1] -> T[1]        |
     |  sel[2] -> T[2]        |
     |  sel[3] -> T[3]        |
     |  sel[4] -> T[4]        |
     |  sel[5] -> T[5]        |
     |  sel[6] -> T[6]        |
     +-------------------------+
           | | | | | | |
           v v v v v v v
         T[6:0] (to analog cells)
\end{verbatim}

\subsection{Control Interface}

\begin{table}[h]
\centering
\begin{tabular}{@{}lcp{7cm}@{}}
\toprule
\textbf{Signal} & \textbf{Width} & \textbf{Description} \\
\midrule
\texttt{sel[0]} & 3 bits & Selects which \texttt{TH[x]} drives \texttt{T[0]} \\
\texttt{sel[1]} & 3 bits & Selects which \texttt{TH[x]} drives \texttt{T[1]} \\
\texttt{sel[2]} & 3 bits & Selects which \texttt{TH[x]} drives \texttt{T[2]} \\
\texttt{sel[3]} & 3 bits & Selects which \texttt{TH[x]} drives \texttt{T[3]} \\
\texttt{sel[4]} & 3 bits & Selects which \texttt{TH[x]} drives \texttt{T[4]} \\
\texttt{sel[5]} & 3 bits & Selects which \texttt{TH[x]} drives \texttt{T[5]} \\
\texttt{sel[6]} & 3 bits & Selects which \texttt{TH[x]} drives \texttt{T[6]} \\
\midrule
\textbf{Total}  & \textbf{21 bits} & Stored in configuration register \\
\bottomrule
\end{tabular}
\caption{Mux Control Signals}
\end{table}

Each \texttt{sel[i]} encodes a value 0--6 selecting the source thermometric bit.

\subsection{Configuration Register}

\begin{itemize}
    \item 21-bit register \texttt{mux\_cfg[20:0]}, partitioned as \texttt{\{sel[6], sel[5], \ldots, sel[0]\}}.
    \item Loaded via SPI/JTAG/APB or hardwired scan chain --- depends on system integration.
    \item \textbf{Default (reset) value:} identity mapping $\rightarrow$ \texttt{sel[i] = i} $\rightarrow$ \texttt{mux\_cfg = 21'b110\_101\_100\_011\_010\_001\_000}.
    \item Register is static during DAC operation; update only when DAC is idle or outputs blanked.
\end{itemize}

\subsection{Implementation}

\begin{lstlisting}[caption={Thermometric Reorder Mux}]
module thermo_reorder_mux (
    input  wire [6:0]  thermo_in,   // from bin2thermo encoder
    input  wire [20:0] mux_cfg,     // {sel[6], sel[5], ..., sel[0]}, 3b each
    output wire [6:0]  thermo_out   // to analog unit cells
);

    // Per-output 7:1 mux
    genvar i;
    generate
        for (i = 0; i < 7; i = i + 1) begin : MUX_STAGE
            wire [2:0] sel = mux_cfg[3*i +: 3];
            assign thermo_out[i] = thermo_in[sel];
        end
    endgenerate

endmodule
\end{lstlisting}

\subsection{Constraints and Validation}

\begin{table}[h]
\centering
\begin{tabular}{@{}lp{10cm}@{}}
\toprule
\textbf{Item} & \textbf{Requirement} \\
\midrule
Uniqueness   & Each \texttt{sel[i]} value must be unique (no two outputs driven by same source). Enforced by configuration software or optional HW checker. \\
Range        & \texttt{sel[i]} $\in \{0\ldots6\}$. Value 7 is illegal; optional: tie output low on illegal select. \\
Glitch-free  & \texttt{mux\_cfg} must not change while DAC is converting. Use clock-domain-safe load (shadow register + sync pulse). \\
Timing       & Adds one 7:1 mux delay ($\sim$1--2 gate delays) to thermometric path. \\
\bottomrule
\end{tabular}
\caption{Mux Constraints}
\end{table}

\subsection{Optional: Hardware Uniqueness Checker}

\begin{lstlisting}[caption={Configuration Uniqueness Checker}]
module uniqueness_check (
    input  wire [20:0] mux_cfg,
    output wire        cfg_valid  // 1 = all selects unique & in range
);

    wire [2:0] sel [0:6];
    wire [6:0] onehot [0:6];
    wire [6:0] combined;

    genvar i;
    generate
        for (i = 0; i < 7; i = i + 1) begin : DECODE
            assign sel[i] = mux_cfg[3*i +: 3];
            assign onehot[i] = (sel[i] < 3'd7) ? (7'b1 << sel[i]) : 7'b0;
        end
    endgenerate

    assign combined = onehot[0] | onehot[1] | onehot[2] | onehot[3]
                    | onehot[4] | onehot[5] | onehot[6];

    wire [6:0] xor_combined = onehot[0] ^ onehot[1] ^ onehot[2] ^ onehot[3]
                            ^ onehot[4] ^ onehot[5] ^ onehot[6];

    assign cfg_valid = (combined == 7'h7F) && (xor_combined == 7'h7F);

endmodule
\end{lstlisting}

%==============================================================================
\section{Top-Level Integration}

\begin{lstlisting}[caption={DAC Digital Top-Level}]
module dac_digital_top (
    input  wire [6:0]  dac_code,     // 7-bit DAC input (MSBs in 2's complement)
    input  wire [20:0] mux_cfg,      // thermometric reorder config
    output wire [3:0]  binary_out,   // to 4 binary-weighted cells
    output wire [6:0]  thermo_out    // to 7 thermometric unit cells
);

    wire [2:0] msb_unsigned;
    wire [6:0] thermo_enc;

    // Binary LSBs pass through
    assign binary_out = dac_code[3:0];

    // 2's complement to unsigned conversion (MSB segment)
    twos_comp_to_unsigned u_tc2u (
        .twos_in      (dac_code[6:4]),
        .unsigned_out (msb_unsigned)
    );

    // MSB unsigned to thermometric conversion
    bin2thermo u_enc (
        .bin    (msb_unsigned),
        .thermo (thermo_enc)
    );

    // Thermometric reorder mux
    thermo_reorder_mux u_mux (
        .thermo_in  (thermo_enc),
        .mux_cfg    (mux_cfg),
        .thermo_out (thermo_out)
    );

endmodule
\end{lstlisting}

%==============================================================================
\section{Verification Checklist}

\begin{table}[h]
\centering
\begin{tabular}{@{}cl@{}}
\toprule
\textbf{\#} & \textbf{Check} \\
\midrule
1 & 2's complement conversion: $-4 \rightarrow 0$, $+3 \rightarrow 7$, monotonic mapping. \\
2 & Monotonicity: consecutive \texttt{dac\_code} increases active elements by exactly one. \\
3 & Full-scale: \texttt{dac\_code[6:4] = 011} ($+3$) $\rightarrow$ all 7 thermo active. \\
4 & Zero-scale: \texttt{dac\_code[6:4] = 100} ($-4$) $\rightarrow$ no thermo elements active. \\
5 & Mid-scale: \texttt{dac\_code[6:4] = 000} ($0$) $\rightarrow$ 4 thermo elements active. \\
6 & Mux identity: default \texttt{mux\_cfg} produces \texttt{thermo\_out == thermo\_enc}. \\
7 & Mux permutation: all $7! = 5040$ valid permutations preserve output set. \\
8 & Invalid config: illegal \texttt{sel} values flagged by \texttt{cfg\_valid}. \\
9 & No glitches on \texttt{thermo\_out} when \texttt{mux\_cfg} is stable. \\
\bottomrule
\end{tabular}
\caption{Verification Checklist}
\end{table}

%==============================================================================
\section{Configuration Examples}

\begin{table}[h]
\centering
\begin{tabular}{@{}llp{6cm}@{}}
\toprule
\textbf{Use Case} & \textbf{mux\_cfg (sel[6]\ldots sel[0])} & \textbf{Effect} \\
\midrule
Identity    & 6, 5, 4, 3, 2, 1, 0 & No reorder \\
Reverse     & 0, 1, 2, 3, 4, 5, 6 & Mirror layout gradient \\
Interleave  & 0, 2, 4, 6, 1, 3, 5 & Spread switching across array \\
Custom cal  & Per-chip calibration  & Compensate measured mismatch \\
\bottomrule
\end{tabular}
\caption{Example Mux Configurations}
\end{table}

%==============================================================================
\section{Deliverables to Analog}

\begin{table}[h]
\centering
\begin{tabular}{@{}llcl@{}}
\toprule
\textbf{Signal} & \textbf{Direction} & \textbf{Width} & \textbf{Description} \\
\midrule
\texttt{binary\_out[3:0]} & Digital $\rightarrow$ Analog & 4 & Binary segment drivers \\
\texttt{thermo\_out[6:0]} & Digital $\rightarrow$ Analog & 7 & Thermometric segment drivers (reordered) \\
\bottomrule
\end{tabular}
\caption{Analog Interface Signals}
\end{table}

%==============================================================================
\section*{Revision History}

\begin{table}[h]
\centering
\begin{tabular}{@{}llll@{}}
\toprule
\textbf{Rev} & \textbf{Date} & \textbf{Author} & \textbf{Description} \\
\midrule
1.0 & 2025-06-01 & --- & Initial release \\
1.1 & \today & --- & Added 2's complement to unsigned conversion stage before thermometric encoder \\
\bottomrule
\end{tabular}
\end{table}

\end{document}
